%%%%%%%%%%%%%%%%%%%%%%%%%%%%%%%%%%%%%%%%%%%%%%%%%%%%%%%%%%%%%%%%%%%%%%%%%%%%%
%  enslyonstage-exemple.tex — FULL EXAMPLE of a completed ENS internship report
%
%  ⚠  This file is for REFERENCE only.
%     To write your own report, open and edit  enslyonstage-gabarit.tex  instead.
%     Copy-paste sections from here as needed.
%
%  Uses the enslyonstage.cls class
%
%  Project structure:
%    enslyonstage-gabarit.tex ← YOUR blank template (start here)
%    enslyonstage-exemple.tex ← this file (full worked example)
%    enslyonstage.cls         ← class file (same directory or on TEXMFHOME)
%    references.bib           ← BibTeX bibliography
%    figures/                 ← images (jpg, pdf, png, eps …)
%    logos/                   ← hautdepage.pdf  basdepage.pdf  Master.jpg
%                                logo_ens.pdf  (+ logo_labo.* optional)
%
%  Compile order:
%    pdflatex enslyonstage-exemple  →  biber enslyonstage-exemple  →  pdflatex
%  On Overleaf: compilation is automatic (biber is used by default).
%
%  Final PDF file name (ENS requirement):
%    LastName-FirstName-ENS-M1.pdf   or   LastName-FirstName-ENS-M2.pdf
%
% ═══════════════════════════════════════════════════════════════════════════
%  ENS DE LYON — OFFICIAL GUIDELINES  (keep these comments during writing)
% ═══════════════════════════════════════════════════════════════════════════
%  FORMAT
%    ✓ Written in English
%    ✓ Two-column layout (this class)        ✓ Pages numbered (automatic)
%    ✓ PDF, preferably < 4–5 MB
%    ✓ File name: "LastName-FirstName - ENS-M1.pdf"  or  "...-M2.pdf"
%    ✗ No supplemental data / annexes
%      (exceptions: contact stages.biosciences@ens-lyon.fr)
%
%  LENGTH & DISPLAY ITEMS
%    M1 : ≤ 5 000 words (± 10 %)   ≤ 6 figures + tables
%    M2 : ≤ 4 000 words (± 10 %)   ≤ 5 figures + tables
%    → word count excludes references and figure/table legends
%    → on Overleaf: Menu → Word Count  (uses the %TC: directives in the class)
%
%  TITLE     ≤ 125 characters including spaces
%  ABSTRACT  150–200 words · English · NO references
%  REFS      ≤ 35 publications · PNAS (new) numbered style [1],[2]…
%            Unpublished data / personal comm.: in-text only, NOT in ref list
%
%  STYLE
%    • Spell out 1–9 unless with units  ("three replicates" / "3 µm", "5 mM")
%    • Units: metric, negative exponents  (µmol m⁻² s⁻¹  NOT  µmol/m²/s)
%    • Gene names: italics  —  Protein names: roman  (follow field conventions)
%    • Abbreviations: define at first use, keep to a minimum
%    • Stats: specify test, n, and exact p-value for every numerical result
% ═══════════════════════════════════════════════════════════════════════════

% ── THEME ──────────────────────────────────────────────────────────────────
% Available (22): orange (default) · violet · blue · red · green · teal · gray
%                 navy · azure · indigo · midnight · slate · purple · pink
%                 crimson · wine · amber · gold · brown · lime · forest · cyan
\documentclass[orange]{enslyonstage}

% ── BIBLIOGRAPHY ───────────────────────────────────────────────────────────
% Bibliography resource (biblatex — compiled with biber, automatic in Overleaf)
% Create a references.bib file with your citations
\addbibresource{references.bib}

% ── CUSTOM SHORTCUTS ───────────────────────────────────────────────────────
% Define frequently-used units as shortcuts to save typing
\newcommand{\um}{\si{\micro\metre}}       % \um   → µm
\newcommand{\uM}{\si{\micro\Molar}}       % \uM   → µM
\newcommand{\degC}{\si{\degreeCelsius}}   % \degC → °C

%%%%%%%%%%%%%%%%%%%%%%%%%%%%%%%%%%%%%%%%%%%%%%%%%%%%%%%%%%%%%%%%%%%%%%%%%%%%%
%  METADATA
%  Fill in your report information below. These commands populate the cover
%  page, headers, footers, and PDF metadata.
%%%%%%%%%%%%%%%%%%%%%%%%%%%%%%%%%%%%%%%%%%%%%%%%%%%%%%%%%%%%%%%%%%%%%%%%%%%%%

% ── TITLE ──────────────────────────────────────────────────────────────────
% ENS requirement: maximum 125 characters including spaces
% Should capture the conceptual significance for a broad audience
%
% TIP — title with a line break: use \texorpdfstring so hyperref metadata
% (PDF title, bookmarks) still receives a plain-text version:
%   \titre{First Part:\texorpdfstring{\\}{ }Second Part of the Title}
%
\titre{Characterisation of a Novel Protein Complex Involved in Post-Transcriptional
	Regulation During Cellular Stress Response}

% ── AUTHOR ─────────────────────────────────────────────────────────────────
% Your name as it should appear on the cover page
% Use \textsc{} for small caps on last name
\auteur{Firstname \textsc{Lastname}}

% ── DEGREE LEVEL ───────────────────────────────────────────────────────────
% Enter either "M1" or "M2" (without "Biosciences" — added automatically)
\master{M1}

% ── BIOLOGICAL DOMAIN ──────────────────────────────────────────────────────
% Shown on cover page and in footer
% Examples: Molecular and Cellular Biology, Neurosciences, Ecology, etc.
\domaine{Molecular and Cellular Biology}

% ── INTERNSHIP DATES ───────────────────────────────────────────────────────
% Format: DD Month YYYY (e.g., 02 February 2025)
\datedebut{02 February 2025}
\datefin{30 June 2025}

% ── LABORATORY ─────────────────────────────────────────────────────────────
% Laboratory name and/or institute
% Example: Lab Name, Institute Name, CNRS UMR XXXX, City
\laboratoire{
    Laboratory of RNA Biology,
    Montpellier, France
    }

% ── SUPERVISORS ────────────────────────────────────────────────────────────
% Lab supervisor (internship mentor) — name and email
\tuteurstage{Dr. Jane Smith}{} % No email in the official
% ENS academic supervisor — name and email
\tuteurENS{Prof. Jean Dupont}{jean.dupont@ens-lyon.fr} % not shown

% ── ABSTRACT ───────────────────────────────────────────────────────────────
% ENS requirements:
%   - 150–200 words (single paragraph)
%   - English language
%   - NO literature references
%   - Include: (1) brief background and question, (2) main results without
%     extensive detail, (3) summary of significance
%   - Keep accessible to non-specialists
\abstracts{%
	Write your abstract here. Single paragraph, 150–200 words, English.
	No literature references permitted. Include: (1)~brief biological background
	and question addressed, (2)~description of main results without extensive
	experimental detail, (3)~summary of significance. Keep language accessible
	to non-specialists.
	\lipsum[1][1-5]
	% ── VISUAL ABSTRACT (optional) ─────────────────────────────────────────
	% Place a graphical abstract (figure, scheme, workflow…) below the text.
	% The optional argument sets the width (default: 0.9\textwidth).
	% Use any image format accepted by \includegraphics.
	\abstractimage[0.5\textwidth]{example-image}
}

% ── KEYWORDS ───────────────────────────────────────────────────────────────
% ENS guideline: 5 to 8 keywords
% Comma-separated list of key terms for your work
\keywords{RNA-binding proteins, stress granules, post-transcriptional regulation,
	proteomics, yeast two-hybrid}

% ── WORD COUNT ─────────────────────────────────────────────────────────────
% ENS limits: M1 ≤ 5000 words (±10%)  /  M2 ≤ 4000 words (±10%)
% Count excludes references and figure/table legends
% On Overleaf: Menu → Word Count  (the %TC: directives in enslyonstage.cls
%              automatically exclude figure captions and table environments)
% Locally:  texcount -sum -1 enslyonstage-exemple.tex
\nombredemots{6 767}

% ── PREVIOUS INTERNSHIPS ───────────────────────────────────────────────────
% Table of previous internships shown on cover page
% Use \stagerow{Year}{Lab, City, Country}{Supervisor}{Topic}
% Leave empty fields for levels not yet completed (e.g., M2 if you're in M1)
\stagesanterieurs{
	\stagerow{L3}{Lab name, City, Country}{Dr.\ Firstname \textsc{Lastname}}{Biological topic}
	\stagerow{M1}{Lab name, City, Country}{Dr.\ Firstname \textsc{Lastname}}{Biological topic}
	\stagerow{M2}{}{}{}
}

%%%%%%%%%%%%%%%%%%%%%%%%%%%%%%%%%%%%%%%%%%%%%%%%%%%%%%%%%%%%%%%%%%%%%%%%%%%%%
%  LOGOS
%  Place logo files in the logos/ directory.
%  Cover page automatically uses hautdepage.pdf and basdepage.pdf.
%%%%%%%%%%%%%%%%%%%%%%%%%%%%%%%%%%%%%%%%%%%%%%%%%%%%%%%%%%%%%%%%%%%%%%%%%%%%%

% ── ENS LOGO ───────────────────────────────────────────────────────────────
% Displayed in footer center at 8mm height (automatic)
% Optional height argument: \logoENS[height]{filepath}
\logoENS[1.5cm]{logos/logo_ens.pdf}

% ── MASTER LOGO ────────────────────────────────────────────────────────────
% Kept for compatibility, not shown in footer if \logoENS is set
\logoformation[2cm]{logos/Master.jpg}

% ── LABORATORY LOGO ────────────────────────────────────────────────────────
% Shown on cover page between domain and info block (optional)
% Remove or comment out if not needed
\logoLabo[3cm]{example-image}

%%%%%%%%%%%%%%%%%%%%%%%%%%%%%%%%%%%%%%%%%%%%%%%%%%%%%%%%%%%%%%%%%%%%%%%%%%%%%
%  HEADER / FOOTER CUSTOMIZATION (optional)
%  Defaults are set in enslyonstage.cls and work well for most reports.
%  Uncomment and modify only if you need custom headers/footers.
%%%%%%%%%%%%%%%%%%%%%%%%%%%%%%%%%%%%%%%%%%%%%%%%%%%%%%%%%%%%%%%%%%%%%%%%%%%%%
% Default layout:
%   top-left  : current section name (\leftmark)
%   top-right : title wrapped at 0.6\linewidth
%   bot-left  : \domaine in small caps
%   bot-centre: \logoENS at 8mm (positioned below footer rule)
%   bot-right : page/total

%\coinhautgauche{custom text}
%\coinhautdroit{custom text}
%\coinbasgauche{custom text}
%\coinbascentre{\includegraphics[height=8mm]{logos/custom.pdf}}
%\coinbasdroit{custom text}

%%%%%%%%%%%%%%%%%%%%%%%%%%%%%%%%%%%%%%%%%%%%%%%%%%%%%%%%%%%%%%%%%%%%%%%%%%%%%
\begin{document}
%%%%%%%%%%%%%%%%%%%%%%%%%%%%%%%%%%%%%%%%%%%%%%%%%%%%%%%%%%%%%%%%%%%%%%%%%%%%%

% ── COVER PAGE ─────────────────────────────────────────────────────────────
% Generates official ENS cover page with all metadata
% Automatically switches to two-column mode after cover
\fairepagedegarde

% ── ABSTRACT BOX ───────────────────────────────────────────────────────────
% Displays abstract in full-width colored box at top of page 1
% Must be called immediately after \fairepagedegarde
\faireabstract

% ── TABLE OF CONTENTS ──────────────────────────────────────────────────────
% Optional — remove if not required by your tutor
% Shows sections, subsections, and subsubsections
\tableofcontents

%%%%%%%%%%%%%%%%%%%%%%%%%%%%%%%%%%%%%%%%%%%%%%%%%%%%%%%%%%%
\section{Introduction}
\label{sec:intro}
%%%%%%%%%%%%%%%%%%%%%%%%%%%%%%%%%%%%%%%%%%%%%%%%%%%%%%%%%%%
% ENS requirements for Introduction:
%   - Provide background information to place subject in general biological
%     context with MORE detail than a standard journal paper
%   - Must be understandable to non-specialists
%   - End with clear statement of the question(s) addressed in the report
%   - Define non-standard abbreviations at first use

% TIP — drop cap on the first letter of the introduction:
%   \lettrine[lines=3]{T}{he} context of this study...
%   The lettrine package is already loaded by enslyonstage.cls.
\paragraph{}
\lipsum[1-4]

% \newpage

%%%%%%%%%%%%%%%%%%%%%%%%%%%%%%%%%%%%%%%%%%%%%%%%%%%%%%%%%%%
\section{Materials and Methods}
\label{sec:methods}
%%%%%%%%%%%%%%%%%%%%%%%%%%%%%%%%%%%%%%%%%%%%%%%%%%%%%%%%%%%
% ENS requirements for Materials and Methods:
%   - Include sufficient detail to allow procedures to be repeated
%   - Standard procedures may be referenced rather than fully described
%   - Statistical tests: specify test name, sample size (n), and exact p-value
%
% Organize with \paragraph{} headings for each experimental approach

\paragraph{Cell culture and strains.}
\lipsum[1][1-4]

\paragraph{Plasmid construction.}
\lipsum[2][1-4]

\paragraph{Protein extraction and western blotting.}
\lipsum[3][1-4]

\paragraph{Co-immunoprecipitation.}
\lipsum[4][1-4]

\paragraph{Mass spectrometry.}
\lipsum[5][1-4]

\paragraph{Statistical analysis.}
% ENS requirement: specify test, sample size n, and exact p-value for each result
\lipsum[6][1-4]

%%%%%%%%%%%%%%%%%%%%%%%%%%%%%%%%%%%%%%%%%%%%%%%%%%%%%%%%%%%
\section{Results}
\label{sec:results}
%%%%%%%%%%%%%%%%%%%%%%%%%%%%%%%%%%%%%%%%%%%%%%%%%%%%%%%%%%%
% ENS requirements for Results:
%   - Present findings clearly and succinctly
%   - Provide concise summary of data in ALL display items (figures and tables)
%   - Subheadings encouraged to organize content
%   - Display item limits: M1 ≤ 6 figures+tables  /  M2 ≤ 5 figures+tables
%
% Figure placement tips:
%   - Include figures as small as possible without compromising clarity
%   - Use \FloatBarrier after each subsection to control figure placement
%   - Legends must be self-sufficient (understandable without reading main text)

\paragraph{}
\lipsum[1]

%-----------------------------------------------------------------------------
%  FIGURE INSERTION — \insererfigure
%  Usage: \insererfigure[options]{filepath}{caption}{label}
%
%  Options (comma-separated, all optional):
%    col     — single-column figure, width=\linewidth
%    p       — full-page float [p], figure alone on its page
%    b       — double-column figure, prefer bottom placement [!hb]
%    h=<dim> — custom image height (default 0.33\textheight)
%    w=<dim> — custom image width (overrides h)
%
%  Default (no options): double-column figure [!htb], h=0.33\textheight
%
%  Examples:
%    \insererfigure{file.png}{Caption text}{fig:label}
%       → double-column, default height
%    \insererfigure[col]{file.png}{Caption}{fig:label}
%       → single-column, width=\linewidth
%    \insererfigure[w=0.6\textwidth]{file.png}{Caption}{fig:label}
%       → double-column, custom width
%    \insererfigure[p,w=0.85\textwidth]{file.png}{Caption}{fig:label}
%       → full-page float
%-----------------------------------------------------------------------------


\subsection{Identification of interaction partners}

% Figure 1 — double column (default, spans both columns)
% ENS guideline: include figures as small as possible without losing clarity
% Legend must be self-sufficient (no need to read the main text to understand it)
\insererfigure[w=0.6\textwidth]%
{example-image}%
{\textbf{Identification of proteins co-purifying with the bait.}
	(\textbf{A})~Schematic of the pull-down strategy.
	(\textbf{B})~Volcano plot ($n = 3$ biological replicates); dashed lines:
	$p < 0.05$ and fold-change $\geq 2$ (log$_2$ scale).
	(\textbf{C})~Validation by co-IP; Input~(I), flow-through~(FT), eluate~(E).
	Scale bar~= 10\,\um. Data are mean~$\pm$~SD.
	See also \autoref{fig:network}.}%
{fig:pulldown}

\lipsum[2-5]

\FloatBarrier

\subsection{Functional validation}

% Figure 2 — single column
% Use [col] option to fit figure within one column
\insererfigure[col]%
{example-image}%
{\textbf{Co-immunoprecipitation confirms the interaction \textit{in vivo}.}
	Input~(I), flow-through~(FT) and eluate~(E) are shown.
	Molecular weight markers (kDa) on the left.
	Three independent experiments gave similar results\cite{Vincent2020}.}%
{fig:coip}

\lipsum[6-7]

\FloatBarrier

% Figure 3 — single column, custom height
\insererfigure[col,h=0.28\textheight]%
{example-image}%
{\textbf{Stress-induced relocalization of the complex.}
	Confocal images of GFP-tagged bait under basal conditions (left) and
	after \SI{30}{\minute} treatment with \SI{1}{\milli\Molar}~\ce{H2O2} (right).
	Scale bar~= \SI{5}{\micro\metre}. $n > 100$ cells per condition,
	$p < 0.001$, Mann–Whitney $U$-test\cite{Journo2021}.}%
{fig:localisation}

\lipsum[8]

\FloatBarrier

\subsection{Structural characterisation}

% Figure 4 — full page
% [p] creates a dedicated float page with the figure alone
% Useful for large, detailed figures
\insererfigure[p,w=0.85\textwidth]%
{example-image}%
{\textbf{Interaction network of the identified complex.}
	Nodes: proteins coloured by functional category.
	Edge thickness: interaction score.
	Outlined nodes: validated by $\geq 2$ independent methods.
	$K_d = \SI{8.5e-9}{\Molar}$ (SPR, $n = 3$)\cite{Trizac2019}.
	Visualised with Cytoscape~3.9.}%
{fig:network}

\lipsum[9-10]

%-----------------------------------------------------------------------------
%  TABLE — full-width (table*) and single-column (table)
%
%  Caption placement: ABOVE the tabular (biology convention)
%  table*  — spans both columns (same as double-column figures)
%  table   — single column
%
%  ENS display item limit: M1 ≤ 6 figures+tables  /  M2 ≤ 5 figures+tables
%    → tables and figures share the same quota
%-----------------------------------------------------------------------------

\begin{table*}[!htb]
  \centering
  \caption{\textbf{Summary of identified interaction partners.}
    Proteins co-purifying with the bait identified by mass spectrometry
    ($n = 3$ biological replicates). Enrichment is expressed as log$_2$
    fold-change (FC) over the control. $p$-values corrected by
    Benjamini–Hochberg. nd, not detected in control.}
  \label{tab:partners}
  \small
  \begin{tabular}{llccc}
    \toprule
    \textbf{Protein} & \textbf{Function} & \textbf{log$_2$ FC} & \textbf{Adj.\ $p$-value} & \textbf{Validation} \\
    \midrule
    ProteinA  & RNA-binding          & 4.2 & \num{1.3e-4} & co-IP, Y2H \\
    ProteinB  & Stress granule       & 3.8 & \num{2.1e-4} & co-IP      \\
    ProteinC  & Translation factor   & 2.9 & \num{8.3e-3} & Y2H        \\
    ProteinD  & mRNA decay           & 2.4 & \num{1.2e-2} & ---         \\
    ProteinE  & Scaffold/unknown     & 1.8 & \num{4.5e-2} & ---         \\
    \bottomrule
  \end{tabular}
\end{table*}

\FloatBarrier

%%%%%%%%%%%%%%%%%%%%%%%%%%%%%%%%%%%%%%%%%%%%%%%%%%%%%%%%%%%
\section{Discussion}
\label{sec:discussion}
%%%%%%%%%%%%%%%%%%%%%%%%%%%%%%%%%%%%%%%%%%%%%%%%%%%%%%%%%%%
% ENS requirements for Discussion:
%   - Explain significance of results (broad biological context)
%   - Critical evaluation of results — do NOT oversell
%   - Discuss shortcomings and how they could be addressed
%   - Propose 2–3 concrete follow-up experiments
%
% This is NOT a journal manuscript — show your capacity to discuss and
% defend your work, including negative results and limitations

\paragraph{}
\lipsum[12-14]

\begin{table}[!htb]
  \centering
  \caption{\textbf{Strains used in this study.}
    All strains are \textit{S.~cerevisiae} BY4741 background unless
    otherwise stated.}
  \label{tab:strains}
  \small
  \begin{tabular}{lcc}
    \toprule
    \textbf{Name} & \textbf{Genotype} & \textbf{Source} \\
    \midrule
    BY4741   & \textit{MATa his3$\Delta$1 leu2$\Delta$0} & \citeauthor{Vincent2020}, \citeyear{Vincent2020} \\
    yFL001   & BY4741 \textit{protA-GFP::HIS3}           & This study         \\
    yFL002   & BY4741 \textit{$\Delta$protB::KanMX}      & This study         \\
    \bottomrule
  \end{tabular}
\end{table}

\TODO{Discuss limitations and propose follow-up experiments.}

% ── LaTeX reference examples (DELETE before submission) ─────────────────────
\subsection*{LaTeX reference examples (delete before submission)}

Below are live examples of every useful construct for a biology report.
Delete this entire subsection before submitting.

\paragraph{Units (siunitx).}

Speed: \SI{30}{\centi\metre\per\second} or \SI{30}{cm.s^{-1}}

Cell diameter: \SI{10}{\micro\metre} (or 10\,\um).

Protein concentration: \SI{200}{\micro\gram\per\milli\litre}.

Reaction temperature: \SI{37}{\degreeCelsius} (or 37\,\degC).

Centrifugation: \SI{1500}{\rpm} for \SI{10}{\minute} at \SI{4}{\degC}.
% Note: \rpm is declared in enslyonstage.cls via \DeclareSIUnit{\rpm}{rpm}

Wavelength: \SI{280}{\nano\metre} or \SI{280}{nm}.

Molarity: \SI{5}{\milli\Molar} NaCl, \SI{2}{\micro\Molar} inhibitor (\uM).

Range: \SIrange{50}{100}{\micro\litre}.

\paragraph{Scientific notation.}

Avogadro: \num{6.022e23}.

Boltzmann: \num{1.38e-23}~J$\cdot$K$^{-1}$.

Dissociation constant: $K_d = \SI{8.5e-9}{\Molar}$.

Fold enrichment: \num{3.2}-fold increase.

\paragraph{Chemical formulae (mhchem).}

\begin{enumerate}
	\item Hydrogen peroxide: \ce{H2O2}
	\item Energy currency: \ce{ATP -> ADP + Pi}
	\item Redox: \ce{NAD+ + H+ + 2e- -> NADH}
	\item Ion: \ce{Ca^2+}, \ce{Mg^2+}
	\item Radioisotope label: \ce{^{32}P}
	\item Buffer: \SI{50}{\milli\Molar} \ce{NaH2PO4}, pH~7.4
\end{enumerate}

\paragraph{Statistics.}

\begin{itemize}
	\item All data are mean~$\pm$~SD ($n = 3$ biological replicates) unless stated
	otherwise. Unpaired two-tailed Student's $t$-test: $p = 0.021$.
	\item One-way ANOVA: $F(2,12) = 8.4$, $p = 0.005$, Tukey post-hoc.
	\item Mann–Whitney $U$-test: $p < 0.001$.
	\item Pearson correlation: $r = 0.95$, $p < 0.001$.
	\item Effect size: $\geq 2$-fold (log$_2$ scale).
	\item $n > 100$ cells per condition.
\end{itemize}

\paragraph{Gene and protein names.}

\begin{description}
	\item[Gene (yeast, italics)] \textit{cdc28}, \textit{act1}.
	\item[Protein (yeast, roman)] Cdc28, Act1.
	\item[Gene (human, italics caps)] \textit{BRCA1}, \textit{TP53}.
	\item[Protein (human, roman caps)] BRCA1, p53.
	\item[Mutant allele] \textit{ade2}$^{\Delta}$.
	\item[Species] \textit{S.~cerevisiae}, \textit{E.~coli}, \textit{C.~elegans}.
\end{description}

\paragraph{Equations.}

Michaelis--Menten:
\begin{equation}
    v = \frac{V_{\max}[S]}{K_m + [S]}
    \label{eq:michaelismenten}
\end{equation}

% TIP — always add \noindent after a displayed equation if you want the
% following sentence to continue the paragraph flush-left (no indent):
\noindent where $V_{\max}$ is the maximal velocity and $K_m$ the substrate
concentration at half-maximal velocity.

Henderson--Hasselbalch:
\begin{equation*}
	\text{pH} = \text{p}K_a + \log\frac{[\text{A}^-]}{[\text{HA}]}
\end{equation*}

Beer--Lambert:
\[
	A = \varepsilon \cdot l \cdot c
\]

Hemoglobin Binding Model (cf. \autoref{eq:michaelismenten}):
\begin{align}
    \theta &= \frac{[L]^n}{K_d^n + [L]^n}
             = \frac{1}{1 + \left(\dfrac{K_d}{[L]}\right)^{\!n}}
    \label{eq:hill} \\[6pt]
    \frac{d\theta}{dt} &= k_{on}[L]^n(1-\theta) - k_{off}\,\theta
    \label{eq:hill-kinetics}
\end{align}

\paragraph{Cross-references.}

% \autoref — single target, coloured (defined by enslyonstage.cls):
See \autoref{fig:pulldown} for the pull-down results,
\autoref{fig:coip}A,B for the co-IP validation,
and \autoref{fig:network} for the full interaction network.
Statistical methods are described in \autoref{sec:methods}.

% \crefcolor — one or more targets, coloured (defined by enslyonstage.cls):
% Useful when citing several figures at once: \crefcolor{fig:coip,fig:network}
% Do NOT add \usepackage{cleveref} — it is already loaded by the class.
See \crefcolor{fig:pulldown,fig:coip} for the pull-down and co-IP data.

\paragraph{Bibliography citations (superscript coloured [n]).}
As shown previously\cite{Vincent2020}, the complex assembles under stress.
These results are consistent with independent reports\cite{Journo2021,Trizac2019}.
For reviews, see\cite{AlazardDany2019,Goudemand2019,Escarguel2015}.

% Citing with author name visible in text:
%   following \citeauthor{Smith2020}\cite{Smith2020} as defined by...
%   → "following Smith [3] as defined by..."

\paragraph{Typography notes.}
\textit{In vivo}, \textit{in vitro}, \textit{in situ} — always italics.

Multiplication sign: $5 \times 10^{6}$ cells (never letter~x).

Plus-minus: $12.3 \pm 0.4$~kDa.

Approximately: $\approx \SI{300}{\nano\metre}$.

Non-breaking space before units and references: \SI{10}{\um} (use \texttt{\textasciitilde}).

En-dash for ranges: 10--20\,\um\ (two hyphens in source).
% ────────────────────────────────────────────────────────────────────────

%%%%%%%%%%%%%%%%%%%%%%%%%%%%%%%%%%%%%%%%%%%%%%%%%%%%%%%%%%%
\section*{Author Contributions}
%%%%%%%%%%%%%%%%%%%%%%%%%%%%%%%%%%%%%%%%%%%%%%%%%%%%%%%%%%%
% ENS requirement (mandatory):
%   - State YOUR personal contribution (which experiments/analyses you performed)
%   - Describe the role of EVERY person who participated in your mentoring
%     or helped in the study
%
% Use initials matching those in author names above.
% Example format: "F.L. designed and performed experiments X and Y..."

\paragraph{}
F.L.\ designed and performed all experiments, analysed the data, and wrote
the manuscript.

J.S.\ conceived the project, provided reagents and supervised the internship.

J.D.\ provided academic supervision and critically revised the manuscript.

%%%%%%%%%%%%%%%%%%%%%%%%%%%%%%%%%%%%%%%%%%%%%%%%%%%%%%%%%%%
\section*{Acknowledgements}
%%%%%%%%%%%%%%%%%%%%%%%%%%%%%%%%%%%%%%%%%%%%%%%%%%%%%%%%%%%
% Optional section to thank people who helped but are not listed in
% Author Contributions (technical staff, facility managers, etc.)

This report was typeset using a custom \LaTeX{} class developed for ENS de Lyon Biosciences internship reports.

%--- Bibliography — ≤ 35 references (ENS guideline) --------------------------
% ENS requirements for bibliography:
%   - Maximum 35 publications
%   - PNAS numbered style: superscript [n] with colored hyperlinks
%   - Range compression: [1-3] instead of [1,2,3]
%   - Unpublished data / personal communications: cite in text only, NOT here
%
% Citation commands:
%   \cite{key}           → superscript [1]
%   \cite{k1,k2,k3}      → superscript [1-3] (auto-compressed)
%
% Balance columns on final page (makes last page look cleaner)
\balance

% Make sure references.bib contains all cited entries
\biblio

% ── SUPPLEMENTARY FIGURES ──────────────────────────────────────────────────
% Place supplementary figures AFTER \biblio.
% Usage: \inserersuppfig[<includegraphics options>]{<file>}{<caption>}{<label>}
% The optional first argument overrides the default width=\textwidth.
% Counter resets to S1, S2, … automatically.
% Cross-reference with \autoref{supp:fig1} → "Supplemental Figure S1"

\inserersuppfig{example-image}{%
  \textbf{Representative confocal micrograph of GFP-tagged Ssa1 in control
  cells.}
  Cells expressing Ssa1-GFP were imaged by confocal microscopy (63× oil,
  0.13\,\um/pixel). Scale bar: \SI{5}{\micro\metre}.
  ($n = 50$ cells from three independent experiments.)%
}{supp:fig1}

\inserersuppfig[width=0.6\textwidth]{example-image}{%
  \textbf{Western blot confirming equal loading.}
  Total protein extract ($\SI{20}{\micro\gram}$ per lane) was separated by
  SDS-PAGE and probed with anti-PGK1 as a loading control.
  Molecular weight markers (kDa) are indicated on the left.%
}{supp:fig2}

\end{document}
